\documentclass[11pt,a4paper]{article}

\usepackage[utf8]{inputenc}
\usepackage[T1]{fontenc}
\usepackage[margin=2.5cm]{geometry}
\usepackage{amsmath,amssymb}
\usepackage{booktabs}
\usepackage{graphicx}
\usepackage[hidelinks]{hyperref}

\title{Attention Mechanisms in Sequence Modelling:\\
       Architectures, Complexity, and Practical Efficiency}
\author{Prepared with LaTeXEditor}
\date{July 2026}

\begin{document}
\maketitle

\begin{abstract}
Attention has become the dominant primitive for sequence modelling, displacing
recurrence in machine translation, language modelling, and increasingly in
vision and the sciences. This paper reviews the mechanism from first
principles, derives its computational and memory profile, and surveys the
families of methods that have been proposed to reduce its quadratic cost:
sparse attention, low-rank and kernel approximations, and IO-aware exact
methods. We argue that the field's framing of ``efficient attention'' as a
single problem obscures a more useful distinction between \emph{training-time}
cost, dominated by the attention matrix, and \emph{inference-time} cost,
dominated by the key--value cache. Methods that address one frequently do
nothing for the other, and several widely cited approximations have been
superseded in practice by exact methods with better memory behaviour. We close
with the evaluation problem: reported speedups are rarely comparable, because
they are measured at different sequence lengths, batch sizes, and levels of
hardware utilisation.
\end{abstract}

\tableofcontents

\section{Introduction}
\label{sec:intro}

For three decades the default approach to modelling sequences was recurrence.
A hidden state was carried forward one step at a time, and the network's
capacity to relate distant positions was mediated entirely by that state's
ability to preserve information across many updates. The long short-term
memory network \cite{hochreiter1997} made this workable by introducing gated
paths through which gradients could flow without vanishing, and for two
decades it and its variants defined the state of the art.

The limitation was structural rather than incidental. Recurrence imposes a
sequential dependency: position $t$ cannot be computed until position $t-1$
has been. On hardware whose advantage is parallelism, this is an expensive
constraint. It also means that the path length between two positions $i$ and
$j$ grows linearly in $|i - j|$, so the number of operations through which a
gradient must travel grows with the distance it must cover.

Attention removes both problems at once. Every position attends directly to
every other, so the path length between any pair is constant, and the
computation for all positions can be performed in parallel. This is the
observation that motivated the Transformer \cite{vaswani2017}, and the
architectural consequences have been far-reaching.

The cost is quadratic. Relating every position to every other requires
$O(n^2)$ pairwise scores for a sequence of length $n$, in both time and
memory. At the sequence lengths of early machine translation work --- a few
hundred tokens --- this was unremarkable. At the lengths now routinely
demanded, tens or hundreds of thousands of tokens --- driven by the shift to
large general-purpose models prompted with substantial context
\cite{brown2020} --- it dominates everything else. A very large literature has grown around reducing it.

This paper has three aims. Section~\ref{sec:background} and
Section~\ref{sec:transformer} establish the mechanism and its cost precisely,
because much of the efficiency literature is imprecise about which term it is
actually reducing. Sections~\ref{sec:sparse} through \ref{sec:io} survey the
main families of approach. Section~\ref{sec:inference} then argues that the
training-time and inference-time problems are distinct, and that conflating
them has caused the field to over-value some methods and under-value others.

\section{Background: from recurrence to alignment}
\label{sec:background}

Attention did not arrive with the Transformer. It was introduced as a fix for
a specific failure of encoder--decoder translation models
\cite{bahdanau2015}. In those models the entire source sentence was compressed
into a single fixed-width vector, and translation quality degraded sharply as
sentences grew longer --- the vector became an information bottleneck.

The proposal was to let the decoder look back at all encoder states rather
than a single summary. At each decoding step $t$ the model computes a
distribution over source positions,
\begin{equation}
  \alpha_{tj} = \frac{\exp(e_{tj})}{\sum_{k=1}^{n} \exp(e_{tk})},
  \qquad
  e_{tj} = a(s_{t-1}, h_j),
  \label{eq:bahdanau}
\end{equation}
where $h_j$ is the encoder state at source position $j$, $s_{t-1}$ the decoder
state, and $a$ a small feed-forward network. The context vector
$c_t = \sum_j \alpha_{tj} h_j$ is then a weighted average of source states,
with the weights chosen by the model.

Two properties of Equation~\ref{eq:bahdanau} carried over into everything that
followed. First, the weights are \emph{content-based}: they depend on what is
in the states, not on their positions. Second, they are normalised, so
attention is a redistribution of a fixed budget rather than an unbounded gain.
Both properties are preserved in the Transformer.

\section{The Transformer}
\label{sec:transformer}

\subsection{Scaled dot-product attention}

The Transformer \cite{vaswani2017} replaces the learned alignment function $a$
with a dot product, which is cheaper and maps directly onto matrix
multiplication hardware. Given queries $Q \in \mathbb{R}^{n \times d_k}$, keys
$K \in \mathbb{R}^{n \times d_k}$ and values $V \in \mathbb{R}^{n \times d_v}$,
\begin{equation}
  \mathrm{Attention}(Q, K, V)
    = \mathrm{softmax}\!\left(\frac{QK^{\top}}{\sqrt{d_k}}\right) V.
  \label{eq:attention}
\end{equation}

The scaling by $\sqrt{d_k}$ is not cosmetic. If the components of $q$ and $k$
are independent with zero mean and unit variance, then $q \cdot k$ has
variance $d_k$. Without the correction, the magnitude of the logits grows with
the head dimension, pushing the softmax into a regime where it is
near-deterministic and its gradient is close to zero. Dividing by $\sqrt{d_k}$
restores unit variance and keeps the function in its responsive range.

\subsection{Multiple heads}

A single attention distribution must commit to one notion of relevance. The
Transformer instead runs $h$ attention operations in parallel on projections
of the input,
\begin{align}
  \mathrm{MultiHead}(Q,K,V) &= \mathrm{Concat}(\mathrm{head}_1, \dots,
    \mathrm{head}_h)\, W^O, \\
  \mathrm{head}_i &= \mathrm{Attention}(QW_i^Q, KW_i^K, VW_i^V).
\end{align}
Because $d_k$ is set to $d_{\text{model}} / h$, the total cost is roughly that
of a single head at full width. The heads are observed empirically to
specialise --- some track syntactic dependencies, others positional offsets,
others appear largely redundant, a finding that motivated later work on head
pruning.

\subsection{Cost}
\label{sec:cost}

Let $n$ be sequence length and $d$ the model dimension. The two terms that
matter are:
\begin{align}
  \text{attention:} \quad & O(n^2 d) \ \text{time}, \quad O(n^2) \ \text{memory} \\
  \text{feed-forward:} \quad & O(n d^2) \ \text{time}, \quad O(nd) \ \text{memory}
\end{align}
The crossover is at $n \approx d$. Below it the feed-forward layers dominate;
above it attention does. For a model with $d = 1024$, sequences shorter than
about a thousand tokens spend most of their time \emph{outside} the attention
operation --- a fact frequently lost in papers reporting large speedups on
short benchmarks.

\begin{table}[htbp]
  \centering
  \caption{Asymptotic cost of attention variants. $n$ is sequence length,
           $d$ model dimension, $w$ window size and $r$ projection rank.}
  \label{tab:complexity}
  \begin{tabular}{llll}
    \toprule
    \textbf{Method} & \textbf{Time} & \textbf{Memory} & \textbf{Exact?} \\
    \midrule
    Full attention \cite{vaswani2017}   & $O(n^2 d)$      & $O(n^2)$   & yes \\
    Sparse \cite{child2019}             & $O(n\sqrt{n}d)$ & $O(n\sqrt{n})$ & no \\
    Longformer \cite{beltagy2020}       & $O(nwd)$        & $O(nw)$    & no \\
    Big Bird \cite{zaheer2020}          & $O(nd)$         & $O(n)$     & no \\
    Linformer \cite{wang2020}           & $O(nrd)$        & $O(nr)$    & no \\
    Performer \cite{choromanski2021}    & $O(nd^2)$       & $O(nd)$    & no \\
    Linear attention \cite{katharopoulos2020} & $O(nd^2)$ & $O(nd)$    & no \\
    FlashAttention \cite{dao2022}       & $O(n^2 d)$      & $O(n)$     & yes \\
    \bottomrule
  \end{tabular}
\end{table}

\section{Sparse attention}
\label{sec:sparse}

The first family restricts which pairs are computed at all. If each query
attends to $k \ll n$ keys, cost falls to $O(nkd)$. The design question is
which $k$.

\paragraph{Fixed patterns.}
Sparse Transformers \cite{child2019} use strided and local patterns chosen so
that any two positions are connected within two attention layers, giving
$O(n\sqrt{n})$ cost. Longformer \cite{beltagy2020} combines a sliding window
with a small number of global tokens attending everywhere --- a design that
suits classification, where a single \texttt{[CLS]} position must see the whole
input.

\paragraph{Theoretical coverage.}
Big Bird \cite{zaheer2020} adds random connections to window and global
attention, and shows the result is a universal approximator of sequence
functions and Turing complete. This is a reassuring result but a weak one: the
construction requires a number of layers that grows with the sequence, so it
does not establish that a fixed-depth model retains full attention's
expressiveness.

\paragraph{The practical difficulty.}
Sparse patterns produce irregular memory access. A theoretical reduction of
$8\times$ in FLOPs frequently yields a much smaller wall-clock gain, because
the resulting kernels utilise the hardware poorly. This gap between arithmetic
and realised performance is the central lesson of the efficiency literature,
and it is what motivated the IO-aware approach of Section~\ref{sec:io}.

\section{Low-rank and kernel methods}
\label{sec:lowrank}

The second family observes that the attention matrix is often approximately
low-rank, and avoids materialising it.

Linformer \cite{wang2020} projects keys and values to a fixed dimension $r$
before the product, giving $O(nr)$ cost. The projection is learned and its
size fixed, which means the model cannot extrapolate to sequences longer than
those it was trained on --- a serious limitation given that long-context
extrapolation is precisely what motivates the work.

Kernel methods take a different route. Observing that softmax attention is
\begin{equation}
  \mathrm{Attention}(Q,K,V)_i
    = \frac{\sum_j \exp(q_i \cdot k_j) v_j}{\sum_j \exp(q_i \cdot k_j)},
\end{equation}
one may replace $\exp(q \cdot k)$ with a kernel $\phi(q) \cdot \phi(k)$ and
reassociate:
\begin{equation}
  \mathrm{Attention}(Q,K,V)_i
    = \frac{\phi(q_i) \sum_j \phi(k_j) v_j^{\top}}
           {\phi(q_i) \sum_j \phi(k_j)}.
  \label{eq:linear}
\end{equation}
The sums are now independent of $i$ and can be computed once, giving $O(n)$
cost \cite{katharopoulos2020}. Performer \cite{choromanski2021} chooses $\phi$
as a random feature map that approximates the exponential kernel in
expectation, with variance controlled by the number of features.

Equation~\ref{eq:linear} also reveals something interesting: in the causal
case the running sums are a recurrent state, so linear attention is a
recurrent network with a particular update rule. The field spent five years
replacing recurrence with attention, and this line of work quietly reintroduces
it --- an observation made explicit in recent state-space models
\cite{gu2023}.

\section{IO-aware exact attention}
\label{sec:io}

The methods above trade accuracy for speed. FlashAttention \cite{dao2022}
argues the trade is often unnecessary, because attention is not
compute-bound at all but \emph{memory-bound}: the expensive operation is
writing the $n \times n$ matrix to high-bandwidth memory and reading it back,
not the arithmetic.

The solution is tiling. Blocks of $Q$, $K$ and $V$ are loaded into on-chip
SRAM, the corresponding block of the output is accumulated using an online
softmax that never requires the full row, and the attention matrix is never
written out. The result is numerically identical to
Equation~\ref{eq:attention} --- it is exact --- while reducing memory from
$O(n^2)$ to $O(n)$ and delivering substantial wall-clock speedups.

This result reframes much of the preceding literature. Approximations that
were justified by the $O(n^2)$ memory term address a cost that an exact method
removes. In practice FlashAttention and its successors have displaced most of
the approximate methods in production systems, and papers proposing new
approximations are now expected to compare against it rather than against
naive attention --- a comparison many of the earlier results would not survive.

\section{Positional information}
\label{sec:positional}

Attention as defined is permutation-equivariant: shuffling the input shuffles
the output identically. Word order must therefore be injected separately.

The original sinusoidal encodings \cite{vaswani2017} add fixed vectors to the
input embeddings. Rotary embeddings \cite{su2021} instead rotate queries and
keys by an angle proportional to position, so that the dot product depends
only on relative offset:
\begin{equation}
  \langle R_{\theta m} q, R_{\theta n} k \rangle = g(q, k, m - n).
\end{equation}
This relative property is the reason RoPE extrapolates better than absolute
schemes, and it has become the default in current models. ALiBi
\cite{press2022} is simpler still --- a linear bias on the attention scores
proportional to distance --- and was explicitly designed for train-short,
test-long generalisation.

\section{The inference problem is a different problem}
\label{sec:inference}

Almost all of the above targets the attention matrix, which is a
\emph{training-time} cost. Autoregressive inference has a different profile.

Generating token $t$ requires attending to all previous positions. Rather than
recompute their keys and values, implementations cache them. The cache has
size
\begin{equation}
  2 \cdot n_{\text{layers}} \cdot n_{\text{heads}} \cdot d_{\text{head}}
  \cdot n_{\text{tokens}} \cdot b,
\end{equation}
with $b$ bytes per element. For a 70B-parameter model at 32{,}000 tokens this
runs to tens of gigabytes --- frequently exceeding the memory taken by the
weights, and bounding how many sequences can be served concurrently.

Crucially, none of the approximations in Sections~\ref{sec:sparse}
and~\ref{sec:lowrank} help with this, because the cache is proportional to
sequence length regardless of which pairs are scored. The methods that do help
attack the head structure instead. Multi-query attention \cite{shazeer2019}
shares a single key--value head across all query heads, dividing the cache by
$n_{\text{heads}}$ at some cost in quality. Grouped-query attention
\cite{ainslie2023} interpolates, sharing across groups, and recovers most of
the quality at most of the saving.

\begin{table}[htbp]
  \centering
  \caption{Key--value cache per token, 32-layer model, $d_{\text{model}} = 4096$,
           16-bit precision.}
  \label{tab:kvcache}
  \begin{tabular}{lrr}
    \toprule
    \textbf{Scheme} & \textbf{KV heads} & \textbf{Cache/token} \\
    \midrule
    Multi-head \cite{vaswani2017}   & 32 & 512\,KB \\
    Grouped-query \cite{ainslie2023} & 8 & 128\,KB \\
    Multi-query \cite{shazeer2019}  & 1  & 16\,KB  \\
    \bottomrule
  \end{tabular}
\end{table}

The distinction matters for how the literature should be read. A paper
reporting an $8\times$ reduction in attention cost may deliver no improvement
whatever to a deployed serving system, if that system's bottleneck is cache
memory rather than score computation.

\section{Self-attention and cross-attention}
\label{sec:cross}

Equation~\ref{eq:attention} is stated for arbitrary $Q$, $K$ and $V$, and the
choice of where those come from determines what the layer does.

In \emph{self-attention} all three derive from the same sequence. The layer
relates positions within one input, and the resulting representation of each
token is contextualised by its neighbours. This is the operation that gives
BERT \cite{devlin2019} its bidirectional character: because no mask is applied,
every token sees every other, and the representation of an ambiguous word is
disambiguated by material on both sides of it.

In \emph{cross-attention} the queries come from one sequence and the keys and
values from another. This is the operation Equation~\ref{eq:bahdanau}
originally described --- a decoder attending to an encoder --- and it remains
the mechanism by which conditioning information enters a generative model,
whether that information is a source sentence, an image, or a retrieved
document.

\paragraph{Masking.}
Autoregressive models require that position $t$ not see positions $> t$, which
is imposed by adding $-\infty$ to the corresponding logits before the softmax:
\begin{equation}
  M_{ij} = \begin{cases}
    0 & j \le i \\
    -\infty & j > i
  \end{cases}
  \qquad
  \mathrm{Attention}(Q,K,V) =
    \mathrm{softmax}\!\left(\frac{QK^{\top}}{\sqrt{d_k}} + M\right)V.
\end{equation}
The mask is what makes the difference between an encoder and a decoder; the
parameters and the operation are otherwise identical. It also halves the useful
work, since the masked half of the matrix is computed and discarded --- one of
the easier savings available to a fused implementation.

\section{What the heads learn}
\label{sec:interpretability}

Because attention weights are normalised and directly inspectable, they invite
interpretation, and a substantial literature has grown around reading them.
The results should be treated with more caution than they usually are.

Empirically, heads in trained models do specialise. Some attend to the
immediately preceding token, some to the syntactic head of the current word,
some to delimiters or to the first token of the sequence. So-called
\emph{induction heads} --- which locate an earlier occurrence of the current
token and attend to whatever followed it --- appear reliably during training
and are implicated in in-context learning.

Two caveats are frequently omitted. First, a large fraction of heads can be
pruned after training with little effect on quality, which suggests the
apparent specialisation of any individual head is not evidence that the model
depends on it. Second, attention weight is not the same as causal influence: a
head may attend strongly to a position whose value vector contributes little
to the output, because the magnitude of $v_j$ matters as much as the weight
$\alpha_{ij}$. Attribution methods that account for both give a different
picture from the attention map alone.

The practical consequence is that attention visualisations are a source of
hypotheses rather than conclusions. They are useful for generating an
explanation, and insufficient for confirming one.

\section{Normalisation and trainability}
\label{sec:norm}

A detail of the architecture that receives less attention than the mechanism
itself, but which determines whether deep models train at all, is where
layer normalisation is placed.

The original formulation applied it after the residual addition:
\begin{equation}
  x_{\ell+1} = \mathrm{LayerNorm}\big(x_\ell + \mathrm{Sublayer}(x_\ell)\big).
  \label{eq:postln}
\end{equation}
This \emph{post-norm} arrangement places the normalisation on the residual
path, so gradients flowing back through many layers are repeatedly rescaled.
In practice deep post-norm Transformers require a learning-rate warmup to
train stably, and become difficult to train at all beyond a few dozen layers.

The \emph{pre-norm} alternative normalises the input to the sublayer instead:
\begin{equation}
  x_{\ell+1} = x_\ell + \mathrm{Sublayer}\big(\mathrm{LayerNorm}(x_\ell)\big).
  \label{eq:preln}
\end{equation}
Here the residual stream is untouched, giving gradients a clean path from the
loss to every layer. Pre-norm models train without warmup and scale to far
greater depth, which is why essentially all current large models use it. The
cost is a small quality penalty at equal depth, generally judged worth paying.

This is a useful example of a broader pattern: a change with no effect on the
function being computed at initialisation can nonetheless determine whether
the model is trainable. Much of the progress in scaling has come from such
changes rather than from new mechanisms.

\section{Beyond text}
\label{sec:vision}

Nothing in Equation~\ref{eq:attention} is specific to language. Given any set
of vectors, attention relates them. The main obstacle to applying it elsewhere
has been the quadratic cost, since the natural tokenisation of other
modalities produces long sequences.

Vision Transformers resolve this by splitting an image into patches --- a
$224 \times 224$ image at $16 \times 16$ patches gives 196 tokens, a
manageable length --- and treating the patches as a sequence. The result
matches or exceeds convolutional networks given sufficient data, though it
underperforms them in the low-data regime, which is usually attributed to the
convolution's built-in assumptions of locality and translation equivariance.
Attention has no such priors and must learn them, which costs data.

The same trade appears wherever attention is applied: it is a more general
mechanism than the architectures it replaces, and generality is paid for in
sample efficiency. Where the inductive bias of a specialised architecture is
correct, that architecture will win at small scale; attention wins when there
is enough data to learn the structure that the bias would have supplied.

\section{Evaluation}
\label{sec:evaluation}

Comparing these methods is harder than it appears, for three reasons.

\paragraph{Length dependence.}
Because attention is quadratic and the feed-forward layers are linear, any
speedup figure is a function of the length at which it was measured. A method
reporting $4\times$ at $n = 16{,}000$ may show no gain at $n = 512$, where
attention is not the bottleneck. Results quoted without the length are
uninterpretable.

\paragraph{Utilisation.}
FLOP reductions and wall-clock reductions diverge, sometimes by an order of
magnitude, because irregular access patterns use the hardware badly. The
correct measure is time, on stated hardware, for a stated implementation.

\paragraph{Task sensitivity.}
Long Range Arena \cite{tay2022} was introduced to standardise comparison, and
remains the most-used benchmark. Its tasks are, however, largely synthetic,
and strong performance on them has not reliably predicted strong performance
on language modelling --- several methods that lead the benchmark are not
competitive on real corpora.

\section{Cost in context: scaling}
\label{sec:scaling}

A final consideration reframes the efficiency question entirely. Empirical
scaling work relates model loss to parameters, data, and compute by power
laws, and shows that for a fixed compute budget there is an optimal allocation
between model size and training tokens. Early large models were substantially
under-trained by this criterion --- too many parameters, too few tokens.

The relevance to attention is that context length is a third axis, and it
interacts with the other two. Extending context is not free even with a linear
attention mechanism, because the tokens must come from somewhere: doubling the
context of every training example either halves the number of examples seen at
fixed compute, or doubles the budget. A method that makes long context
\emph{possible} does not make it \emph{worthwhile}, and whether the additional
context earns its cost is a question about the data rather than the
architecture.

This is visible in practice. Models advertised with very long context windows
frequently show degraded retrieval accuracy in the middle of that window ---
the so-called ``lost in the middle'' effect --- because the training
distribution contained few examples requiring genuine use of distant context.
The mechanism supports the length; the training did not teach the model to use
it. Efficiency work that reduces the cost of long sequences addresses only
half of the problem, and arguably the easier half.

We therefore suggest that the useful question is not ``how cheaply can we
attend over $n$ tokens'' but ``what is the marginal value of the $n$th token,
and does it exceed the marginal cost''. The first question is architectural
and largely solved by the methods of Section~\ref{sec:io}; the second is
empirical, dataset-dependent, and substantially less well studied.

\section{Open problems}
\label{sec:open}

Three questions seem to us underdetermined by the current literature.

\paragraph{Is the softmax necessary?}
Equation~\ref{eq:attention} normalises scores into a distribution, which
guarantees the output is a convex combination of values and bounds its
magnitude. Kernel methods (Section~\ref{sec:lowrank}) abandon the exact
softmax and lose little, which raises the question of what the normalisation
buys. One answer is competition: because weights sum to one, attending to one
position costs attention elsewhere, and this pressure may be what drives head
specialisation. If so, unnormalised variants should show measurably less
specialisation, which as far as we are aware has not been tested directly.

\paragraph{How much context is used?}
Reported context windows describe capacity, not behaviour. Measuring how much
of an available window a model actually exploits requires interventional
experiments --- perturbing distant context and measuring the effect on the
prediction --- rather than benchmark accuracy. The few such studies that exist
suggest effective context is substantially shorter than nominal context, but
the measurement is not standardised and results are not comparable across
models.

\paragraph{What replaces attention?}
State-space models \cite{gu2023} achieve linear scaling with a recurrent
formulation and are competitive at moderate scale. Whether they match
attention at frontier scale, and on tasks requiring precise retrieval from
distant context, remains genuinely open. The honest summary is that attention
is currently dominant for empirical rather than theoretical reasons, and a
mechanism with better asymptotics that matched its quality would displace it.
Several candidates exist; none has yet done so decisively.

\section{Conclusion}
\label{sec:conclusion}

Attention replaced recurrence because it removed a sequential dependency and
shortened gradient paths, at the price of quadratic cost. The subsequent
literature has attacked that cost from three directions, and their standing
today is uneven. Sparse and low-rank approximations are theoretically
attractive but frequently disappointing in wall-clock terms, and are largely
superseded by exact IO-aware methods that were possible all along and simply
required attention to the memory hierarchy rather than the operation count.

The more durable lesson is the one in Section~\ref{sec:inference}: ``efficient
attention'' names two different problems. Training cost is dominated by the
attention matrix; inference cost is dominated by the key--value cache. Methods
that address one may do nothing for the other, and a field that reports a
single efficiency number obscures this. We would encourage authors to state
plainly which regime a result applies to, at what sequence length, and against
which baseline.

\begin{thebibliography}{99}

\bibitem{ainslie2023}
J.~Ainslie, J.~Lee-Thorp, M.~de~Jong, Y.~Zemlyanskiy, F.~Lebrón, and
S.~Sanghai.
\newblock GQA: Training generalized multi-query transformer models from
  multi-head checkpoints.
\newblock In \emph{EMNLP}, 2023.

\bibitem{bahdanau2015}
D.~Bahdanau, K.~Cho, and Y.~Bengio.
\newblock Neural machine translation by jointly learning to align and
  translate.
\newblock In \emph{ICLR}, 2015.

\bibitem{beltagy2020}
I.~Beltagy, M.~E. Peters, and A.~Cohan.
\newblock Longformer: The long-document transformer.
\newblock \emph{arXiv:2004.05150}, 2020.

\bibitem{brown2020}
T.~Brown et~al.
\newblock Language models are few-shot learners.
\newblock In \emph{NeurIPS}, 2020.

\bibitem{child2019}
R.~Child, S.~Gray, A.~Radford, and I.~Sutskever.
\newblock Generating long sequences with sparse transformers.
\newblock \emph{arXiv:1904.10509}, 2019.

\bibitem{choromanski2021}
K.~Choromanski et~al.
\newblock Rethinking attention with performers.
\newblock In \emph{ICLR}, 2021.

\bibitem{dao2022}
T.~Dao, D.~Y. Fu, S.~Ermon, A.~Rudra, and C.~Ré.
\newblock FlashAttention: Fast and memory-efficient exact attention with
  IO-awareness.
\newblock In \emph{NeurIPS}, 2022.

\bibitem{devlin2019}
J.~Devlin, M.-W. Chang, K.~Lee, and K.~Toutanova.
\newblock BERT: Pre-training of deep bidirectional transformers for language
  understanding.
\newblock In \emph{NAACL-HLT}, 2019.

\bibitem{gu2023}
A.~Gu and T.~Dao.
\newblock Mamba: Linear-time sequence modeling with selective state spaces.
\newblock \emph{arXiv:2312.00752}, 2023.

\bibitem{hochreiter1997}
S.~Hochreiter and J.~Schmidhuber.
\newblock Long short-term memory.
\newblock \emph{Neural Computation}, 9(8):1735--1780, 1997.

\bibitem{katharopoulos2020}
A.~Katharopoulos, A.~Vyas, N.~Pappas, and F.~Fleuret.
\newblock Transformers are RNNs: Fast autoregressive transformers with linear
  attention.
\newblock In \emph{ICML}, 2020.

\bibitem{press2022}
O.~Press, N.~A. Smith, and M.~Lewis.
\newblock Train short, test long: Attention with linear biases enables input
  length extrapolation.
\newblock In \emph{ICLR}, 2022.

\bibitem{shazeer2019}
N.~Shazeer.
\newblock Fast transformer decoding: One write-head is all you need.
\newblock \emph{arXiv:1911.02150}, 2019.

\bibitem{su2021}
J.~Su, Y.~Lu, S.~Pan, B.~Wen, and Y.~Liu.
\newblock RoFormer: Enhanced transformer with rotary position embedding.
\newblock \emph{arXiv:2104.09864}, 2021.

\bibitem{tay2022}
Y.~Tay, M.~Dehghani, S.~Abnar, Y.~Shen, D.~Bahri, P.~Pham, J.~Rao, L.~Yang,
  S.~Ruder, and D.~Metzler.
\newblock Long Range Arena: A benchmark for efficient transformers.
\newblock In \emph{ICLR}, 2021.

\bibitem{vaswani2017}
A.~Vaswani, N.~Shazeer, N.~Parmar, J.~Uszkoreit, L.~Jones, A.~N. Gomez,
  Ł.~Kaiser, and I.~Polosukhin.
\newblock Attention is all you need.
\newblock In \emph{NeurIPS}, 2017.

\bibitem{wang2020}
S.~Wang, B.~Z. Li, M.~Khabsa, H.~Fang, and H.~Ma.
\newblock Linformer: Self-attention with linear complexity.
\newblock \emph{arXiv:2006.04768}, 2020.

\bibitem{zaheer2020}
M.~Zaheer et~al.
\newblock Big Bird: Transformers for longer sequences.
\newblock In \emph{NeurIPS}, 2020.

\end{thebibliography}

\end{document}
